% ============================================================
% Chapter 3: Project Management
%   Section: Deviations from the Initial Plan
% Included from chapters/03-project-management/project-management.tex
% ============================================================

\section{Deviations from the Initial Plan}
\label{sec:deviations}

The plan was rebalanced three times, always at 900 hours
(Table~\ref{tab:rebalance}).

\begin{table}[H]
\centering
\small
\caption{Hours per phase after each replanning (docs/roadmap.md). The
earlier columns are regrouped into the current phases so that the
three add up the same way.}
\label{tab:rebalance}
\begin{tabular}{@{}lrrr@{}}
\toprule
Phase & Sep 28 & Sep 29 & Sep 30 \\
\midrule
0: Foundations                      & 300 & 300 & 300 \\
1: Pipeline completion              & 140 & 140 & 140 \\
2: Development                      & 140 & 130 & 120 \\
3: Neo4j and infrastructure         &  70 &  65 &  50 \\
4: Evaluation                       &  40 &  80 &  90 \\
5: AI model testing                 & 120 &  80 &  80 \\
6: Interface                        &  40 &  20 &  35 \\
Testing and QA, transversal         &   0 &  50 &  50 \\
7: Thesis report                    &  50 &  35 &  35 \\
\midrule
\textbf{Total}                      & \textbf{900} & \textbf{900} & \textbf{900} \\
\bottomrule
\end{tabular}
\end{table}

\paragraph{28 September: the validation set is a task of its own.} The
custom validation set had been expected to fit into hours freed in
Sprint~3. It turned out to be about 40 hours of work by itself ---
finding articles, extracting claims, searching evidence and labelling
150 facts --- and it was given those 40 hours, paid for by the writing
benchmark (Sprint~6, 70 to 50 hours) and the interface (Sprint~8, 60 to
40).

\paragraph{29 September: evaluation and testing get hours.} Three
things every later result depends on had no hours: the validation set's
calendar, a harness that runs the pipeline over the labelled sets, and
testing as a continuous activity. The validation set got its own
window (29~September to 30~October, about ten facts a day, then the
blind re-label a week later); testing became a transversal track of 50
hours, weighted towards the end; the harness got 40 hours; and the
search fix, without which a benchmark would measure the search engine's
uptime, got 10 hours inside Sprint~4. They were paid for by cutting the
user-facing application to a read-only reader view (accounts, profiles,
bookmarks and recommendations answer no research question), by
reducing the writing benchmark to one local and one hosted model, and
by de-scoping the topic classifier's accuracy, which is on no research
question's path.

\paragraph{30 September: dates follow from hours.} The previous plan
kept the hours honest but not the calendar: one week held 65 hours of
infrastructure while the 20-hour interface had two. Every block's dates
were recomputed from its hours at the 30--35 hours a week the first two
phases had run at; development and evaluation became two tracks run
side by side; and the Neo4j integration, done on 29~September in far
less than its 35 hours, gave 15 hours to the interface.

\paragraph{5 October: experiments before the verdict.} Eighteen hours
of experiments on discovery, selection, admission and the topic
classifier were added without funding; how they are expected to be
absorbed is in Section~\ref{subsec:phases}.

\paragraph{After the plan.} The 30~September plan was overtaken in both
directions.

\begin{itemize}
    \item \textbf{Work done ahead of its dates.} The production compose
    file and the cloud infrastructure (planned for 20~October --
    9~November) were done on 1--2~October, except the first real
    deployment, which needs a billed project and is left undone; the
    reader view (planned for 11--17~December) on 2--4~October; the
    evaluation harness, its metrics and retrieval evaluation (planned
    for 20~October -- 9~November) on 2--4~October; continuous
    integration and the full-pipeline integration test on 2~October;
    and the X-Fact pilot on 13~October.
    \item \textbf{Work added outside the plan.} Discovery by topic
    group, the choice of up to twenty candidates by a person or the AI
    and the newsroom flow (5~October); the mission screen and the
    candidate ranking (6~October); the alignment of the rounds with the
    mission and the groups (9~October); and the person's topic emphasis
    (13~October). Each came from what a live round showed --- election
    coverage in a Culture round, a match report ranked above an
    analysis of decarbonisation --- and each was measured before it was
    kept (Section~\ref{sec:design-selection}).
    \item \textbf{Work behind its dates.} The labelling: by
    14~October the set held 109 natural facts, of which 31 were
    labelled by the annotator and 78 by an AI labelling pass not yet
    reviewed (Subsection~\ref{subsec:custom-procedure}); 39 facts from
    17 off-mission articles had been moved out of the set; and no
    natural claim was \texttt{FALSE}.
\end{itemize}

\paragraph{14 October: feature freeze and thesis priorities.} With
about a third of this document written, mostly on scraping, and no
result yet on the questions the thesis asks, the plan was reordered by
what the thesis needs rather than by the calendar. The author fixed the
four research questions; nothing new is built unless a priority needs
it; and the remaining work was sorted into three levels
(Table~\ref{tab:priorities}). Writing news from verified sources moved
to the lowest level, since no research question asks for it.

\begin{table}[ht]
\centering
\small
\caption{Thesis priorities from 14 October 2026 (docs/roadmap.md).
Where they disagree with the calendar of
Section~\ref{sec:planning}, they win.}
\label{tab:priorities}
\begin{tabularx}{\linewidth}{@{}lX@{}}
\toprule
Level & Items \\
\midrule
P0 --- without these there is no thesis & Fix the research questions
and their metrics (done, 14~October); human adjudication of the gold
set and the agreement between the AI pass and the annotator; a
\texttt{FALSE} class (done, 14~October); a frozen tuning/test split;
the verification runs on the test split under three evidence
conditions; stage attribution; writing, about 140 hours, alongside
everything else \\
\addlinespace
P1 --- what makes it a good thesis & Retrieval as an ablation on frozen
claims; evidence against the age of the claim; two more verifier
models; ten blind and ten assisted selection rounds (RQ3); the
discovery numbers on five days (RQ1) \\
\addlinespace
P2 --- only if time is left & Tuning the ranking and confidence weights
and the pertinence gate; the remaining experiments before the verdict;
the writing-model benchmark; writing news from verified sources \\
\bottomrule
\end{tabularx}
\end{table}
